\documentclass[handout]{beamer}
%\documentclass{beamer}

%\setbeamercovered{dynamic}

\usetheme{default}
\usecolortheme{rose}
\useinnertheme{rounded}
%\useinnertheme[shadow]{rounded}

%\usetheme{Boadilla} 
\usefonttheme{serif}



%\usetheme{Darmstadt}
%\usefonttheme[onlylarge]{structurebold}
\setbeamerfont*{frametitle}{size=\normalsize,series=\bfseries}
\setbeamertemplate{navigation symbols}{}
\setbeamertemplate{footline}[frame number]

%\setbeamercovered{transparent}



\usepackage[english]{babel}
\usepackage[latin1]{inputenc}
\usepackage{times}
\usepackage[T1]{fontenc}
\usepackage{amsmath,amsthm,amssymb}
\usepackage{graphicx}
\usepackage{mathrsfs}
%\usepackage{caption}
%\usepackage{subcaption}
\usepackage{color}


\AtBeginSubsection[]{

  \frame<beamer>{

    \frametitle{Outline}

    \tableofcontents[currentsection,currentsubsection]

  }

}

\def\P{{\cal P}}
\def\si{\bar{\sigma}}
\def\E{{\cal E}}
\def\LL{{\cal L}}
\def\FF{{\cal F}}
\def\SS{{\cal S}}
\def\C{{\cal C}}
\def\H{{\cal H}}
\def\X{{\cal X}}
\def\Y{{\cal Y}}
\def\A{{\cal A}}
\def\B{{\cal B}}
\def\vX{\boldsymbol{X}}
\def\vY{\boldsymbol{Y}}
\def\vZ{\boldsymbol{Z}}
\def\O{{\cal O}}
\def\M{{\cal M}}
\def\I{{\cal I}}
\def\eps{{\varepsilon}}
\def\ch{{\mbox{\rm ch}\hspace{0.4mm}}}
\def\sh{{\mbox{\rm sh}}}
\def\th{{\mbox{\rm th}}}
\def\Av{{\mbox{\rm{Av}}}}
\def\n{\noindent}
\def\bbe{\vskip .15pc\hangindent=.5pc\hangafter=1}
\def\la{{\Bigl\langle}}
\def\ra{{\Bigr\rangle}}
\def\goin{\to\infty}
\def\qed{\hfill\break\rightline{$\bbox$}}

\newcommand{\ow}{\overline{w}}
\newcommand{\ou}{\overline{u}}
\newcommand{\ov}{\overline{v}}
\newcommand{\e}{\mathbb{E}}
\newcommand{\p}{\mathbb{P}}
\newcommand{\Reals}{\mathbb{R}}
\newcommand{\Natural}{\mathbb{N}}
\newcommand{\calA}{\mathcal{A}}
\newcommand{\F}{{\cal F}}
\newcommand{\veps}{\vec{\varepsilon}}
\newcommand{\vsi}{\boldsymbol{\sigma}}
\newcommand{\vrho}{{\boldsymbol{\rho}}}
\newcommand{\vbeta}{\vec{\beta}}
\newcommand{\vtau}{\boldsymbol{\tau}}
\newcommand{\veta}{\vec{\beta}}
\newcommand{\vx}{\mathbf{x}}
\newcommand{\vy}{\mathbf{y}}
\newcommand{\vz}{\mathbf{z}}
\newcommand{\vw}{\boldsymbol{w}}
\newcommand{\vW}{\boldsymbol{W}}
\newcommand{\Exval}[1]{\mathbb{E}#1}
\newcommand{\prob}[1]{\mathbb{P}(#1)}
\newcommand{\Prob}[1]{\mathbb{P}\{#1\}}
\newcommand{\bigProb}[1]{\mathbb{P}\Bigl\{#1\Bigr\}}
\newcommand{\bigprob}[1]{\mathbb{P}\left(#1\right)}
\newcommand{\showcite}[1]{\ensuremath{\stackrel{\textnormal{#1}}{\textnormal{\cite{#1}}}}}
\newcommand{\indi}{\ensuremath{\boldsymbol 1}}
\newcommand{\bB}{\mathbf{B}}
\newcommand{\bW}{\mathbf{W}}
\newcommand{\bu}{\mathbf{u}}
\newcommand{\bv}{\mathbf{v}}
\newcommand{\bF}{\mathcal{F}}
\newcommand{\bC}{\mathcal{C}}
\newcommand{\bL}{\mathcal{L}}
\newcommand{\bX}{\mathcal{X}}
\newcommand{\bY}{\mathcal{Y}}
\newcommand{\bZ}{\mathcal{Z}}

\newcommand{\Pp}{\mathcal{P}}
\newcommand{\s}{\boldsymbol{\sigma}}
\newcommand{\ii}{\iota}
\newcommand{\bt}{\boldsymbol{t}}

\newtheorem{proposition}{\bf Proposition}

\title{Multiple levels of replica symmetry breaking}
\author{Antonio Auffinger\\
	Northwestern University
	\\\vspace{1cm} Based on joint works with\\
	 Wei-Kuo Chen (U. Minnesota), and Qiang Zeng (Northwestern)}

%\institute{University of Minnesota}


\date{}
\begin{document}

\frame{\titlepage}

%\section[Outline]{}


\begin{frame}\frametitle{Parisi measures}
	\noindent Given a (random) function 
\begin{color}{blue}	
	$$ \H: \Sigma \to \mathbb R $$
\end{color}	
and $E \in \mathbb R$, what can we say about the level sets
\begin{color}{blue}	
$$ \mathcal L (E) := \{ \sigma \in \Sigma : \H(\sigma) = E  \}? $$
\end{color}
\end{frame}

\begin{frame}\frametitle{A classic problem}
	\noindent Given a (random) function 
\begin{color}{blue}	
	$$ \H: \Sigma \to \mathbb R $$
\end{color}	
and $E \in \mathbb R$, what can we say about the level sets
\begin{color}{blue}	
$$ \mathcal L (E) := \{ \sigma \in \Sigma : \H(\sigma) = E  \}? $$
\end{color}
\end{frame}


	
%	\pause
%	
%	\begin{figure}
%		\centering
%		\begin{subfigure}[b]{0.3\textwidth}
%			\includegraphics[width=\textwidth]{assignment.pdf}
%		\end{subfigure}\hspace{0.3in}
%	%	~ %add desired spacing between images, e. g. ~, \quad, \qquad, \hfill etc. 
%		%(or a blank line to force the subfigure onto a new line)
%		\begin{subfigure}[b]{0.22\textwidth}
%			\includegraphics[width=\textwidth]{people.pdf}
%		\end{subfigure}
%	%	~ %add desired spacing between images, e. g. ~, \quad, \qquad, \hfill etc. 
%		%(or a blank line to force the subfigure onto a new line)
%	\end{figure}
%	
%	\pause
%	
%%	\begin{figure}[h!]
%%		\centering
%%		\includegraphics[width=0.45\textwidth]{score.pdf}
%%		%  \caption{SK model}
%%	\end{figure}
%	
%	\pause 
%	
%	Dean's problem: Find the optimizer of
%	\begin{align*}
%	\max_{\sigma\in\{-1,+1\}^N}\sum_{i,j=1}^Ng_{ij}\sigma_i\sigma_j.


\begin{frame}\frametitle{A classic problem}
	\noindent 
\begin{color}{blue}	
	$$ \H: \Sigma \to \mathbb R, \quad \mathcal L (E) :=   \{ \sigma \in \Sigma : \H(\sigma) = E \pm \varepsilon  \}. $$
\end{color}

\vspace{0.1cm}
Some ways to address the question: 
	
\begin{itemize}
\item $\Sigma = \Sigma_{N}$ a metric space,  $\H = \H_{N}$; $\varepsilon = \varepsilon_{N}$ small.
\item Choose two points in the level set $\mathcal L (E)$ and say something about their distance. 
\item Topology of the level set. Is $\mathcal L(E)$ connected? How many connected components? How many holes?
\end{itemize}

\end{frame}


\begin{frame}{A picture}
\begin{figure}
\begin{center}
    \includegraphics[scale=0.30]{energylandscape.jpg} 
\end{center}
\end{figure}
\end{frame}
%\begin{frame}
%	\frametitle{Edwards-Anderson model}
	
%	\begin{itemize}
%		\item Consider a finite graph $(V,E)$ on $\mathbb{Z}^d.$ 
%		\item Hamiltonian (or Energy): For $\sigma\in \{-1,1\}^V,$
%		$$
%		H(\sigma)=\sum_{(i,j)\in E}g_{ij}\sigma_i\sigma_j,
%		$$
%		where $g_{ij}$ are i.i.d. $N(0,1).$
%		\item How to compute $
%		\max_{\sigma\in\{-1,1\}^N}H(\sigma)?$
%		\item Frustration:
%		\begin{figure}[h!]
%			\centering
%			\includegraphics[width=0.3\textwidth]{frustration.pdf}
%			\caption{Frustration}
%		\end{figure}
%	\end{itemize}
	
%\end{frame}

\begin{frame}{A thermodynamical approach}

	%The limiting free energy associated to the (inverse) temperature $\beta\in\mathbb{R},$ is defined as 
%	$$F(\beta):=\lim_{N\rightarrow\infty} \frac{1}{N} \log Z_{N}(\beta),$$
	Fix $\beta>0$, the Gibbs measure associated to $\H_{N}$ is defined as
    		\begin{align*}
    		G_{N,\beta}(d\sigma)&=\frac{\exp\bigl(\beta \H_N(\sigma)\bigr)d\sigma}{Z_N(\beta)},
    		\end{align*}
where $$Z_N(\beta):=\int_{\Sigma}\exp \big( \beta \H_N(\sigma)\big)d\sigma.$$ 

    		Denote by $\left<\cdot\right>_\beta$ the Gibbs average with respect to the Gibbs measure $G_{N,\beta}$ 
		\begin{align*}
    		F_N(\beta)&=\frac{1}{N}\log Z_{N}(\beta).
    		\end{align*}
\end{frame}



\begin{frame}
	\frametitle{Equivalence of ensembles}
	
There exists a map between temperature parameter and energy levels: 	
	$$\begin{color}{blue} \beta \mapsto E \end{color} $$
	
	
Let $\beta>0$ be fixed. Assume for some $\delta >0$  there exists a nonrandom function $F:(\beta-\delta,\beta+\delta)\rightarrow\mathbb{R}$ such that for any $\beta'\in(\beta-\delta,\beta+\delta)$, 
$$    		\lim_{N\rightarrow\infty}F_N(\beta')= \lim_{N\rightarrow\infty}\frac{1}{N}\log Z_{N}(\beta')=F(\beta'), \,\,a.s.$$

\begin{theorem} If $F$ is differentiable at $\beta,$ then
  $$ \lim_{N\rightarrow\infty}  \la  \Big |\frac{\mathcal H_N(\sigma)}{N}-F'(\beta) \Big |  \ra_\beta=0 \quad a.s.$$
  \end{theorem}
  
 \end{frame}


% % % % % % % % % % % % % % % % % % % % % % % % % % % % % % % % % % % % % % % % % % % % % % % % % % % % % % % % % % % % % % % % % % % % % % % %



\begin{frame}
	\frametitle{Equivalence of ensembles}
Under differentiability of the limiting free energy:	
	$$ \begin{color}{blue} E(\beta) = F'(\beta) \end{color}$$
\begin{itemize}	
\item Gibbs measure concentrates at the level set $E(\beta) = F'(\beta).$
\vspace{0.1cm}
\item If $F$ is not differentiable, the Hamiltonian stays inside the interval given by the left and right derivatives at $\beta$.
\vspace{0.1cm}
\item One tries to understand the geometry of the level set $E(\beta)$ by understanding the Gibbs measure at inverse temperature $\beta$. 
\end{itemize}

\end{frame}


\begin{frame}{Thermodynamical approach}
	
	\pause
\begin{itemize}
	\item Sample $\sigma$ and $\sigma'$ independently from the Gibbs measure $G_{N}$. Goal: understand the limiting distribution of the distance 
	$$ d(\sigma_{1}, \sigma_{2}) \text{ under } G_{N}^{2}.$$
	
	\vspace{0.1cm}
	\item Fix $A \subset \mathbb R$, look at the constrained coupled free energy: 
	$$ \int_{\Sigma \times \Sigma} \mathbf 1_{\{d(\sigma_{1}, \sigma_{2}) \in A\}} e^{\beta (\mathcal H(\sigma_{1}) + \mathcal H(\sigma_{2}))} d \sigma_{1} d \sigma_{2}.  $$

	
	\end{itemize}
		
\end{frame}

\begin{frame}{Example 1: O'Connell-Yor semi-discrete polymer model}
	
	\pause
\begin{itemize}
	\item Consider the configuration space $\Sigma_{N}= \{ \sigma \in \mathbb R^{N+1} : 0 = \sigma_{0} \leq \ldots \leq \sigma_{N-1} \leq \sigma_{N} = N.\}$ Set 
	
$$ \mathcal H(\sigma) = \sum_{i=1}^{N} (B_{i}(\sigma_{i}) - B_{i}(\sigma_{i-1})),$$ 
where $(B_{i})_{1\leq i \leq N}$ are i.i.d. standard Brownian motions.	 
	\vspace{0.1cm}
	\item Free energy explicit given by:
	
	$$\begin{color}{blue}F(\beta)=
- \mathcal L [-\psi](-\beta^{2})-\log(\beta^{2}), \quad \beta\neq 0, \end{color}$$ 
where $\psi$ is the digamma function, and $\mathcal L[-\psi]$ is the Legendre transform of $-\psi$:
$$\mathcal L[-\psi] (x) = \sup_{\theta>0} [x\theta + \psi(\theta)].$$ 
	\end{itemize}
		
\end{frame}

\begin{frame}{Example 1: O'Connell-Yor semi-discrete polymer model}

\begin{itemize}
\item O'Connell-Yor, Moriarty-O'Connell, Seppalainen-Valko, Borodin-Corwin-Ferrari, Borodin-Corwin, ... 
\item Explicit derivative in this case: 
	\begin{align*}
	\begin{color}{blue}F'(\beta)=2 \bigg(\beta \psi_{1}^{-1}(\beta^{2}) - \frac{1}{\beta}\bigg). \end{color}
 	\end{align*}
\item Differentiability also expected  for other polymer models.
\end{itemize}

\end{frame}

\begin{frame}
	\frametitle{Example 2: Spin glasses}

\pause 

\begin{itemize}
	\item Spin Glasses are alloys with strange magnetic properties. Ex: CuMn
	
	\pause
	    
	    \begin{itemize}
	    	\item[-] Physically: Spin + Glass
	  
	        \item[-] Mathematically: Randomness + Frustration
	
		\end{itemize}
		
	\pause
	
	\item Spin glasses exhibit many crucial features commonly shared in many real world problems:
	\begin{itemize}
		\item[-] Traveling salesman problem.
		\item[-] Hopfield neural network.
	\end{itemize}

\end{itemize}

\end{frame}




\begin{frame}{Example 2: The Sherrington-Kirkpatrick (SK) model}

\begin{itemize}
	\item The SK model is defined by the Hamiltonian or energy
	
	\begin{align*}
	\mathcal H_N(\sigma)=\frac{1}{\sqrt{N}}\sum_{i,j=1}^Ng_{ij}\sigma_i\sigma_j\qquad\Bigl(=\frac{1}{\sqrt{N}}\left(g\sigma,\sigma\right)\Bigr)
	\end{align*}
    for each $\sigma\in\{-1,+1\}^N$, where
	$g_{ij}$ are i.i.d. $N(0,1)$.
	
	\item {\color{red}Frustration}:  
	
		\begin{figure}[h!]
			\centering
			\includegraphics[width=0.20\textwidth]{frustration.pdf}
			%  \caption{SK model}
		\end{figure}		
\end{itemize}
	

\end{frame}

\begin{frame}{Example 2: The Sherrington-Kirkpatrick (SK) model}

\begin{itemize}
	\item The SK model is defined by the Hamiltonian or energy
	
	\begin{align*}
	\mathcal \H_N(\sigma)=\frac{1}{\sqrt{N}}\sum_{i,j=1}^Ng_{ij}\sigma_i\sigma_j\qquad\Bigl(=\frac{1}{\sqrt{N}}\left(g\sigma,\sigma\right)\Bigr)
	\end{align*}
    for each $\sigma\in\{-1,+1\}^N$, where
	$g_{ij}$ are i.i.d. $N(0,1)$.
			
\item It is a Gaussian process with
	 	$$
	 	\e \H_N(\sigma^1)\H_N(\sigma^2)=N(R(\sigma^1,\sigma^2)\bigr)^2,
	 	$$
	 	where
	 	$$R(\sigma^1,\sigma^2)=\frac{1}{N}\sum_{i=1}^N\sigma_i^1\sigma_i^2.$$
	 				
\end{itemize}
	

\end{frame}




\begin{frame}{Example 2: The mixed$p$-spin}

\begin{itemize}
\item Mixed $p$-spin model: A Gaussian process with
	 	$$
	 	\e \H_N(\sigma^1)\H_N(\sigma^2)=N \xi\bigg(R(\sigma^1,\sigma^2)\bigg),
	 	$$
	 	where
	 	$$R(\sigma^1,\sigma^2)=\frac{1}{N}\sum_{i=1}^N\sigma_i^1\sigma_i^2.$$
	 				\item SK: $\xi(t)=t^{2}$. 
				  	\item Pure $p$-spin: $\xi(t)=t^{p}$.
\end{itemize}

\end{frame}




\begin{frame}{The limit free energy - Parisi Formula}
\begin{theorem}[M. Talagrand '06, Panchenko '10  - Parisi Formula]
For all $\beta > 0$ the following limit exists a.s.:
\[
\lim_{N \rightarrow \infty} \frac{1}{N} \log Z_N = F(\beta)
\]
\end{theorem}
\begin{itemize}
\item $F(\beta)=F(\beta, \xi)$ is given by a variational principle that for large $\beta$ is highly untrivial:
\begin{equation*}
\begin{color}{blue}F(\beta) = \inf_{\rho \in M[0,1]} \mathcal P(\beta,\rho). \end{color}
\end{equation*}
\item Predicted by G. Parisi. 
\item \begin{color}{red}  What is the minimizer? What is the role of the minimizer? \end{color}
\end{itemize}
\end{frame}


\begin{frame}{The minimizer of the Parisi Functional}


\textbf{Parisi measure.}

\vspace{0.5cm}
We call a measure that minimizes the previous variational problem a {\it Parisi measure}. We denote it by $\mu_P = \mu_{P}(\beta)$. \\
\vspace{0.2cm}


\begin{itemize}

\item {\rm (Guerra '03)} Minimizer exists.
\item {\rm (A.-Chen '14)} Minimizer is unique.

\vspace{0.2cm}
It is expected that:
\vspace{0.2cm}

\item Up to symmetry, $\mu_P$ is the limiting law of the overlap under $\mathbb E G_N^{2}$.

\item The limit law of overlap fully describes the measure $G_N$ as $N$ goes to infinity.
\end{itemize}

\end{frame}


\begin{frame}{Predictions}
\includegraphics[scale=0.30]{slide13}
\end{frame}




\begin{frame}{More predictions}
\textbf{For any mixed $p$-spin, any Parisi measure should satisfy the following:} 
\vspace{0.5cm}
\begin{itemize}
\item[$(1)$]  The origin is contained in the support of the Parisi measure at any temperature.
\vspace{0.5cm}
\item[$(2)$]  One expects FRSB behavior at low temperature.
\vspace{0.5cm}
\item[$(3)$]  Any Parisi measure has a jump discontinuity at the top of the support at any temperature.
\end{itemize}
\end{frame}





\begin{frame}
	\frametitle{Significance of the Parisi measure}
    
    \noindent Two major predictions:
    
    \begin{enumerate}
    	\item[(1)] Full Replica Symmetry Breaking:
    	\begin{figure}[h!]
    		\centering
    		\includegraphics[width=0.7\textwidth]{SK-known.pdf}
    		%  \caption{SK model}
    	\end{figure}
    	
    	\begin{theorem}[A.-Chen-Zeng '17] The cardinality of $\mbox{supp} \; {\color{red}\mu_\beta}$ diverges as $\beta\rightarrow \infty.$ (FRSB holds at zero temperature).
    	\end{theorem}
    \end{enumerate}
\end{frame}



\begin{frame}
\begin{enumerate}
		\item[(2)] Ultrametricity: with probab.$\approx 1$, for i.i.d. $\sigma^1,\sigma^2,\sigma^3\thicksim G_N$,
		 $$\|\sigma^1-\sigma^2\|{\color{red}\leq} \max\bigl(\|\sigma^1-\sigma^3\|,\|\sigma^2-\sigma^3\|\bigr).$$
		\begin{figure}[h!]
			\centering
			\includegraphics[width=0.7\textwidth]{tree.pdf}
			%  \caption{SK model}
		\end{figure}
		
		
		Panchenko '11: Ultrametricity for generic mixed $p$-spin models.
\end{enumerate}


\end{frame}
%
%
%\begin{frame}\frametitle{Parisi formula for the maximal energy}
%
%	\begin{theorem}[A.-Chen'16]
%	
%	\begin{itemize}
%		\item Parisi formula at zero temperature:
%		{\footnotesize$$
%			\lim_{N\rightarrow\infty}\e\max_{\sigma\in\{-1,+1\}^N}\Bigl(\frac{\mathcal H_N(\sigma)}{N}\Bigr)=\inf_{\gamma}\Bigl(\Psi_{\gamma}(0,h)-\frac{1}{2}\int_0^1s\gamma(s)ds\Bigr)
%			$$	
%	    }
%		\item Minimizer $\gamma_P$ exists and is unique. Minimizer describes the landscape near the ground state.
%	\end{itemize}
%	\end{theorem}
%	
%%	{\color{red}Random optimization}$\Rightarrow$ {\color{red}Nonrandom convex optimization}
%
%\end{frame}

%     
\begin{frame}{Energy landscape}
	\begin{itemize}

		\item Can we say anything about the energy landscape of $\mathcal H_N$?

		\item In the bulk (REM Universality): Bovier, Ben Arous-Kuptsov.   
		
		\item Near the global maximum/minumum: landscape is highly non-trivial.
		
		%\item We focus on the mixed even $p$-spin model.
	
    \end{itemize}
\end{frame}
 



\begin{frame}{Spherical model}
	
	 \begin{itemize}
	 	\item {\color{red}Spherical} mixed $p$-spin model, 
	 	\begin{align*}
	 	\H_N(\sigma)&=\sum_{p\in \mathbb{N}}c_p^{1/2}\Bigl(\frac{1}{N^{(p-1)/2}}\sum_{1\leq i_1,\ldots,i_p\leq N}g_{i_1,\ldots,i_p}\sigma_{i_1}\cdots\sigma_{i_p}\Bigr)
	 	\end{align*}
	 	for $\sigma\in S^{N-1}:=\{\sigma\in\mathbb{R}^N:\sum_{i=1}^N\sigma_i^2=N\}.$
	 	\pause
	 	
	 	\item
	 	Note
	 	$$
	 	\e \H_N(\sigma^1)\H_N(\sigma^2)=N\xi\bigl(R(\sigma^1,\sigma^2)\bigr),
	 	$$
	 	where $$\xi(x)=\sum_{p\in \mathbb{N}}c_ps^p\,\,\mbox{for}\,\,\sum_{p\in \mathbb{N}}c_p=1.$$ 
	  
	 \end{itemize}
\end{frame}

\begin{frame}{Example: Spherical SK model}
	\begin{itemize}
		\item Write
	\begin{align*}
	\H_N(\sigma)=\frac{1}{\sqrt{N}}\sum_{i,j=1}^Ng_{ij}\sigma_i\sigma_j=\frac{1}{\sqrt{N}}\left(g\sigma,\sigma\right)=\frac{1}{\sqrt{2N}}\left(W\sigma,\sigma\right),
	\end{align*}
	where $W$ is a GOE matrix, i.e., $W_{ij}=W_{ji}$, $\e W_{ij}=0$, $\e W_{ij}^2=1.$
	
	\pause
	
	\item Then $$\max_{\sigma\in S^{N-1}}\H_N(\sigma)=\frac{\mbox{Largest eigenvalue of $W$}}{\sqrt{2N}}.$$
	\item The critical points of $H_N$ are $\pm$eigenvectors of $W.$
	
	\item Easy to describe the landscape: 2 local minima, 2 saddle points of any index.  
	
\end{itemize}

\end{frame}

\begin{frame}{Energy landscape: near the minimum}
\begin{figure}
\begin{center}
    \includegraphics[scale=0.30]{energylandscape.jpg} 
\end{center}
\end{figure}
\end{frame}



	
\begin{frame}{Energy landscape: multiple peaks}
	 Overlap $R(\sigma,\sigma')=\frac{1}{N}\sum_{i=1}^N\sigma_i\sigma_i'.$ 
	\pause
	
	\begin{theorem}[Multiple peaks, A.-Chen '17 ]
	For any $\varepsilon>0,$ there exists a constant $K>0$ s.t. for any $N\geq 1,$ with probability at least $1-Ke^{-N/K},$ $\exists S_N \subset \Sigma_{N}$ such that 
		\begin{itemize}
			\item[$(i)$] {\color{red}$|S_N|\geq e^{N/K}$},
			\item[$(ii)$]  $\forall\sigma\in S_N$,
			$
			\Bigl|\frac{\H_N(\sigma)}{N}-\min_{\sigma'\in\Sigma_N}\frac{\H_N(\sigma')}{N}\Bigr|< \varepsilon.
			$
			%\item[$(iii)$] $\forall\sigma\in S_N,$ $|m(\sigma)|<\varepsilon.$
			\item[$(iv)$]  $\forall\sigma,\sigma'\in S_N$ with $\sigma\neq \sigma'$, $|R(\sigma,\sigma')|< \varepsilon$.
		\end{itemize}
	\end{theorem}
	
	\pause
	
		\begin{itemize}
			\item (1) Chatterjee '09: $|S_N|\geq (\log N)^c$ 
			\item (2) Ding-Eldan-Zhai '14: $|S_N|\geq N^c.$
			\item Chen-Handschy-Lerman '18: $\Sigma_{N}=\{\pm 1 \}^{N}.$
			\item (1) and (2): larger class of Gaussian Processes
		\end{itemize}
\end{frame}







\begin{frame}{Sphere: Crisanti-Sommers formula}

	\begin{theorem}[Chen-Sen '15, Jagannath-Tobasco '16]
	\begin{align*}
			GSE&:=\lim_{N\rightarrow\infty}\min_{\sigma\in S_N}\Bigr(\frac{\H_N(\sigma)}{N}+h\frac{\sum_{i=1}^{N}\sigma_i}{N}\Bigr)\\
			&=\inf_{\nu}\frac{1}{2}\Bigl(\int_0^1(\xi'(s)+h^2)\nu(ds)+\int_0^1\frac{dq}{\nu((q,1])}\Bigr),
			\end{align*}
			where the infimum is over all $\nu(ds)=\gamma(s)1_{[0,1)}(s)ds+\Delta \delta_1(ds).$
	\end{theorem}

\begin{itemize}
	\item Minimizer exists and is unique. We call $\nu_P$ the Parisi measure.

	\item Denote $$
	\nu_P(ds)=1_{[0,1)}(s)\gamma_P(s)ds+\Delta_P\delta_1(s).$$
	
	\pause
	
	\item What does $\gamma_P$ look like? 
		\begin{itemize}
			\item[-] {\color{red}Replica Symmetry (RS):} {\color{blue}$\gamma_P\equiv 0$}
			\item[-] {\color{red}One Replica Symmetry Breaking (1RSB):} {\color{blue}$\gamma_P\equiv C>0$}
			\item[-] {\color{red}Full Replica Symmetry Breaking (FRSB):} {\color{blue}$\gamma_P(ds)=f(s)ds$}
        \end{itemize}
\end{itemize}

\end{frame}


	\begin{frame}{1 RSB scenario}
		
		\begin{theorem}
				Assume $h=0$ and $\xi(s)=s^p$ for even $p\geq 3$. For any $\varepsilon>0$, there exist $\eta,K>0$ such that for all $N\geq 1,$
				\begin{align}\label{1RSB:eq1}
				\p\bigl(\exists \; \mbox{$\sigma^1,\sigma^2\in \mathcal{L}_N(\eta)$ s.t. $|R(\sigma^1,\sigma^2)|\in [\varepsilon,1-\varepsilon]$}\bigr)&\leq Ke^{-\frac{N}{K}}.
				\end{align}
		\end{theorem}
	
		
			\begin{figure}[h!]
				\centering
				\includegraphics[width=1\textwidth]{1RSB.pdf}
				%  \caption{SK model}
			\end{figure}
			
	$\mathcal{L}_N(\eta)$ is given by a union of exponentially many nearly orthogonal components. 
	     	
		
	\end{frame}

	
	\begin{frame}{1 RSB phase}
	
	
	
		\begin{itemize}
		\item No path of low energy connecting different ``components''. Existence of diverging barriers.
		\item Number of connected components is exponentially large in $N$: same applies to the Euler characteristic of the level set. 
		
		
		\item A.-Ben Arous-Cerny '13 + Subag '15 $\Rightarrow$ similar results for pure $p$-spin by counting the number of critical points.
		\item Theorem above holds for a large collection of $\xi$'s,  e.g., 
			 $$\xi(s)=cs^p+(1-c)s^{p+2}\,\,\mbox{for even $p\geq 4.$}$$
			$$\xi(s)=\frac{e^{s^2}-1-s^2}{e-2}.$$
		\end{itemize}
	\end{frame}


\begin{frame}
	\frametitle{Energy landscape: Full Replica Symmetry Breaking}
	\begin{theorem} If $\xi'(1)+h^2<\xi''(1)$ and $\xi''(s)^{-1/2}$ is concave on $(0,1]$, then
		{\footnotesize$$
			\gamma_0(s)=
			\left\{
			\begin{array}{ll}
			0,&\mbox{if $q\in[0,q_0)$},\\
			{\color{red}\frac{\xi'''(s)}{2\xi''(s)^{3/2}}},&\mbox{if $q\in [q_0,1)$}
			\end{array}
			\right.
			$$}
		and
		{\footnotesize$$
			GSE=q_0\xi''(q_0)^{1/2}+\int_{q_0}^1\xi''(q)^{1/2}dq,
			$$}
		where $q_0\in[0,1]$ is the unique solution to $\xi'(q_0)+h^2=q_0\xi''(q_0).$ 
	\end{theorem}

\end{frame}

\begin{frame}{Energy landscape: Full Replica Symmetry Breaking}

	\begin{theorem}[A. - Chen '17]
		Let $h\neq 0$. Assume $\xi'(1)+h^2<\xi''(1)$ and $\frac{1}{\sqrt{\xi''}}$ is concave on $(0,1]$.
			\begin{enumerate}
				\item Let $u\in [q_0,1].$ For any $\varepsilon,\eta>0,$ there exists $K>0$ such that
				\begin{align*}
				%\label{FRSB:eq2}
				\p\bigl(\exists \; \mbox{$\sigma^1,\sigma^2\in \mathcal{L}_N(\eta)$ s.t. {\color{red}$|R(\sigma^1,\sigma^2)-u|<\varepsilon$}}\bigr)&\geq 1-Ke^{-\frac{N}{K}}.
				\end{align*}
				\item For any $\varepsilon>0$, there exist $\eta,K>0$ such that
				\begin{align*}
				%\label{FRSB:eq3}
				\p\bigl(\exists \; \mbox{$\sigma^1,\sigma^2\in \mathcal{L}_N(\eta)$ s.t. {\color{red}$R(\sigma^1,\sigma^2)\in [-1,q_0-\varepsilon]$}}\bigr)&\leq Ke^{-\frac{N}{K}}.
				\end{align*}
			\end{enumerate}
	\end{theorem}
	
		\begin{figure}[h!]
			\centering
			\includegraphics[width=0.9\textwidth]{FRSB.pdf}
		\end{figure}
\end{frame}

\begin{frame}{Main ingredients of the proof}
\begin{itemize}
\item Guerra-Talagrand's bound: control of the constrained partition function.

\item Stochastic optimization: handle the behavior of the solution of  Parisi's formula.

\item Connect the support of the Parisi measure to the possible distances between points in the level set. 
\end{itemize}
\end{frame}



\begin{frame}\frametitle{Open questions}
	
	    \pause
	    
	\begin{itemize}
		\item SK model:
	    \begin{itemize}
                    \item[-] Quantitative dependence of the energy landscape on $N.$
                    \item[-] Further structure of the energy landscape.
   		        	\item[-] Fluctuation.
                    %\item[-] Number of spin configurations $\sigma\in \{-1,1\}^N$ s.t. 
                    %\begin{align*}
                    %\sigma_i&=\mbox{sign}\Bigl(\frac{1}{\sqrt{N}}\sum_{j=1}^Ng_{ij}\sigma_j\Bigr),\,\,\forall i=1,\ldots,N.
                    %\end{align*}
		\end{itemize}
		
		\pause
		
		\item Spherical model:
		 \begin{itemize}
                    \item[-]  A.- Zeng '18: Examples of 2-RSB models. 
                    \item[-]  Determine the phase diagram for the energy landscape.
                    %\item[-] Number of spin configurations $\sigma\in \{-1,1\}^N$ s.t. 
                    %\begin{align*}
                    %\sigma_i&=\mbox{sign}\Bigl(\frac{1}{\sqrt{N}}\sum_{j=1}^Ng_{ij}\sigma_j\Bigr),\,\,\forall i=1,\ldots,N.
                    %\end{align*}
		\end{itemize}
		
		\item Bipartite SK model: Let $N_1=cN$ and $N_2=(1-c)N.$
		{\footnotesize\begin{align*}
		H_N(\sigma)=\frac{1}{\sqrt{N}}\sum_{i=1}^{N_1}\sum_{j=1}^{N_2}g_{ij}\tau_i\rho_j
		\end{align*}}
		for  $\sigma=(\tau,\rho)\in \{-1,+1\}^{N_1}\times \{-1,+1\}^{N_2}.$ Note
		{\footnotesize\begin{align*}
		\e H_N(\sigma)H_N(\sigma')=c(1-c)NR(\tau,\tau')R(\rho,\rho').
		\end{align*}}
		Questions:
		\begin{itemize}
			\item[-] Free energy?
			\item[-] Maximum energy?
			\item[-] Energy landscape?
		\end{itemize}
	
		
		\end{itemize}
	
	\end{frame}
     

\begin{frame}
	\begin{center}
		\Large \bf Thank you for your attention.
	\end{center}
\end{frame}
\end{document}

% % % % % % % % % % % % % % % % % % % % % % % % % % % % % % % % % % % % % % % % % % % % % % % % % % % % % % % % % % % % % % % % % % % % % % % % % % % % % %
% % % % % % % % % % % % % % % % % % % % % % % % % % % % % % % % % % % % % % % % % % % % % % % % % % % % % % % % % % % % % % % % % % % % % % % % % % % % % % % % % % % % % % % % % % % % % % % % % % % % % % % % % % % % % % % % % % % % % % % % % % % % % % % % % % % % % % % % % % % % % % % % % % % % % % % % % % %
